\documentclass[aspectratio=169]{beamer}
\usetheme{Madrid}
\usepackage[spanish,es-nodecimaldot]{babel}
\usepackage[utf8]{inputenc}
\usepackage[T1]{fontenc}
\usepackage{amsmath}
\usepackage{amssymb}
\usepackage{graphicx}
\usepackage{booktabs}
\usepackage{xcolor}

% Color UNS
\definecolor{azulUNS}{RGB}{30,60,120}
\usecolortheme[named=azulUNS]{structure}

\title[MAD1 -- U6]{Métodos de Análisis de Datos I}
\subtitle{Unidad 6 -- Bloque 3: Comunicación de resultados}
\author{Departamento de Matemática -- UNS}
\date{2026}

\begin{document}

\begin{frame}
  \titlepage
\end{frame}

\begin{frame}{Contenidos}
  \tableofcontents
\end{frame}

%-----------------------------------------------
\section{El análisis no termina en el modelo}
%-----------------------------------------------

\begin{frame}{El análisis no termina en el modelo}
  Un modelo entrenado, con buenas métricas, es el \textbf{inicio} de la comunicación, no el final.

  \bigskip
  \begin{columns}
    \column{0.5\textwidth}
      \textbf{Lo que el analista sabe:}
      \begin{itemize}
        \item Qué modelo se usó y por qué
        \item Qué métricas se optimizaron
        \item Qué supuestos se hicieron
        \item Qué limitaciones tiene
      \end{itemize}
    \column{0.5\textwidth}
      \textbf{Lo que el destinatario necesita:}
      \begin{itemize}
        \item Qué significa el resultado para su problema
        \item Cuánto puede confiar en él
        \item Qué hacer con esa información
        \item Qué no puede hacerse con ese análisis
      \end{itemize}
  \end{columns}

  \bigskip
  \begin{alertblock}{La brecha}
    La distancia entre ``el modelo tiene AUC = 0.87'' y ``deberían cambiar la política de stock'' es comunicación. Es trabajo del analista cerrar esa brecha.
  \end{alertblock}
\end{frame}

%-----------------------------------------------
\section{Conocer a la audiencia}
%-----------------------------------------------

\begin{frame}{Conocer a la audiencia}
  No existe una sola forma de comunicar resultados. Depende de quién escucha.

  \bigskip
  \begin{columns}
    \column{0.33\textwidth}
      \textbf{Técnica}\\
      (colegas, revisores)
      \begin{itemize}
        \item Mostrar el pipeline completo
        \item Métricas detalladas
        \item Incertidumbre cuantificada
        \item Código reproducible
      \end{itemize}
    \column{0.33\textwidth}
      \textbf{Tomadores de decisión}\\
      (gerentes, directores)
      \begin{itemize}
        \item Resultado en términos del negocio
        \item Confianza, no fórmulas
        \item Implicancias y riesgos
        \item Recomendación concreta
      \end{itemize}
    \column{0.33\textwidth}
      \textbf{Público general}\\
      (comunicación, prensa)
      \begin{itemize}
        \item Analogías y ejemplos
        \item Sin jerga técnica
        \item Una idea central
        \item Visual dominante
      \end{itemize}
  \end{columns}
\end{frame}

\begin{frame}{El error más común: hablarle a uno cuando estás con el otro}
  \begin{exampleblock}{Situación real}
    Un analista presenta ante el directorio de una empresa de logística. Muestra la curva ROC, explica el umbral de clasificación, detalla la matriz de confusión.\\
    \bigskip
    La directora pregunta: \textit{``¿Y esto en plata qué significa?''}
  \end{exampleblock}

  \bigskip
  \textbf{Lo que debería haberse dicho:}\\
  \textit{``Con este modelo, de cada 100 envíos que hoy se retrasan sin que nadie lo anticipe, podríamos detectar 78 con un día de anticipación. Eso equivale a reducir un 30\% los costos de reenvío urgente.''}

  \bigskip
  Traducir métricas técnicas a consecuencias concretas es una habilidad tan importante como saber construir el modelo.
\end{frame}

%-----------------------------------------------
\section{Tipos de hallazgos y cómo presentarlos}
%-----------------------------------------------

\begin{frame}{Tipo 1: confirmás lo que ya sospechaban}
  \begin{columns}
    \column{0.55\textwidth}
      \textbf{La situación:}\\
      El análisis confirma algo que los dueños del problema ya intuían o sabían por experiencia.

      \bigskip
      \textbf{Reacción típica:}\\
      \textit{``Eso ya lo sabíamos. ¿Para qué sirvió el análisis?''}

      \bigskip
      \textbf{Cómo encuadrarlo:}
      \begin{itemize}
        \item Cuantificar lo que antes era intuición.
        \item Mostrar la magnitud, no solo la dirección.
        \item Establecer una línea de base objetiva.
        \item ``Ahora lo sabemos con evidencia.''
      \end{itemize}
    \column{0.45\textwidth}
      \begin{exampleblock}{Ejemplo}
        Análisis de ventas de una cadena de supermercados: ``Las ventas caen los lunes.''\\
        \bigskip
        La respuesta: \textit{``Obvio, siempre fue así.''}\\
        \bigskip
        El valor real: \textit{``Caen un 23\% en promedio, con mayor impacto en la sección de frescos. Eso justifica ajustar el pedido del viernes en esa categoría específica.''}
      \end{exampleblock}
  \end{columns}
\end{frame}

\begin{frame}{Tipo 2: descubrís algo contraintuitivo}
  \begin{columns}
    \column{0.55\textwidth}
      \textbf{La situación:}\\
      El análisis arroja un resultado que contradice la creencia establecida o la experiencia del área.

      \bigskip
      \textbf{Reacción típica:}\\
      \textit{``El modelo debe estar mal.''} o \textit{``Eso no puede ser.''}

      \bigskip
      \textbf{Cómo manejarlo:}
      \begin{itemize}
        \item No defender el modelo, explorar juntos.
        \item Verificar si hay error en los datos o en el pipeline.
        \item Si el resultado se sostiene, presentar la evidencia paso a paso.
        \item Separar \textit{``el modelo dice X''} de \textit{``por lo tanto deben hacer X''}.
      \end{itemize}
    \column{0.45\textwidth}
      \begin{exampleblock}{Ejemplo}
        Un análisis de RR.HH. muestra que los empleados con mayor capacitación tienen mayor rotación.\\
        \bigskip
        La intuición: capacitar retiene. Los datos: los capacitados son más empleables y se van.\\
        \bigskip
        El análisis no dice que no capaciten: dice que la retención requiere otras palancas además de la capacitación.
      \end{exampleblock}
  \end{columns}
\end{frame}

\begin{frame}{Tipo 3: los datos no alcanzan para responder la pregunta}
  \begin{columns}
    \column{0.55\textwidth}
      \textbf{La situación:}\\
      El análisis no puede responder la pregunta original: datos insuficientes, variable objetivo no disponible, sesgo de selección, etc.

      \bigskip
      \textbf{Reacción típica (incorrecta):}\\
      Forzar una respuesta, minimizar las limitaciones, o presentar resultados débiles como si fueran robustos.

      \bigskip
      \textbf{Lo correcto:}
      \begin{itemize}
        \item Decirlo claramente.
        \item Explicar \textit{qué} falta y \textit{por qué} importa.
        \item Proponer cómo obtener esos datos en el futuro.
        \item Ofrecer lo que sí se puede responder.
      \end{itemize}
    \column{0.45\textwidth}
      \begin{alertblock}{Regla de oro}
        Un analista que dice ``con estos datos no se puede concluir eso'' es más valioso que uno que presenta una conclusión sin fundamento.\\
        \bigskip
        La credibilidad a largo plazo se construye siendo honesto sobre los límites del análisis.
      \end{alertblock}
  \end{columns}
\end{frame}

%-----------------------------------------------
\section{Visualización efectiva}
%-----------------------------------------------

\begin{frame}{El gráfico correcto para el mensaje correcto}
  \begin{columns}
    \column{0.5\textwidth}
      \textbf{Quiero mostrar\ldots}
      \begin{itemize}
        \item Distribución de una variable $\rightarrow$ histograma, boxplot, violín
        \item Relación entre dos variables $\rightarrow$ dispersión, correlación
        \item Evolución en el tiempo $\rightarrow$ línea
        \item Comparación entre categorías $\rightarrow$ barras
        \item Proporción del todo $\rightarrow$ barras apiladas (no torta)
        \item Importancia relativa $\rightarrow$ barras horizontales ordenadas
        \item Incertidumbre $\rightarrow$ barras de error, intervalos de confianza
      \end{itemize}
    \column{0.5\textwidth}
      \begin{alertblock}{El gráfico de torta}
        Los humanos somos malos comparando ángulos. Una barra es casi siempre mejor que una torta.\\
        Excepción válida: mostrar que algo es ``más de la mitad'' o ``casi todo''.
      \end{alertblock}
      \bigskip
      \begin{exampleblock}{Regla práctica}
        Si el gráfico requiere una leyenda para entenderse, preguntarse si una tabla sería más clara.
      \end{exampleblock}
  \end{columns}
\end{frame}

\begin{frame}{Errores frecuentes en visualización}
  \begin{enumerate}
    \item \textbf{Eje Y que no empieza en cero} (en barras): exagera diferencias visualmente.
    \item \textbf{Escala inapropiada}: lineal cuando debería ser logarítmica, o viceversa.
    \item \textbf{Demasiada información en un solo gráfico}: el mensaje se pierde.
    \item \textbf{Colores sin semántica}: usar rojo/verde para cosas que no son malas/buenas.
    \item \textbf{Títulos descriptivos en lugar de argumentativos}: ``Ventas por mes'' vs. ``Las ventas cayeron consistentemente en el segundo trimestre''.
    \item \textbf{Omitir la incertidumbre}: mostrar la media sin el intervalo de confianza puede ser engañoso.
    \item \textbf{Confundir correlación con causalidad} en el título o la narrativa.
  \end{enumerate}
\end{frame}

\begin{frame}{El título como argumento}
  \begin{columns}
    \column{0.5\textwidth}
      \textbf{Título descriptivo:}\\
      ``Relación entre edad y salario''

      \bigskip
      \textbf{Título argumentativo:}\\
      ``El salario crece con la experiencia, pero se estanca después de los 45 años''

      \bigskip
      El título argumentativo:
      \begin{itemize}
        \item Dice qué ver en el gráfico.
        \item Reduce la ambigüedad de interpretación.
        \item Hace al gráfico autosuficiente.
        \item Obliga al analista a tener una tesis.
      \end{itemize}
    \column{0.5\textwidth}
      \begin{alertblock}{Atención}
        Un título argumentativo implica responsabilidad. Si el gráfico no sostiene la tesis, el título es deshonesto.\\
        \bigskip
        Antes de escribir un título argumentativo: verificar que el gráfico realmente lo muestre.
      \end{alertblock}
  \end{columns}
\end{frame}

%-----------------------------------------------
\section{Narrativa y estructura del reporte}
%-----------------------------------------------

\begin{frame}{Estructura de un reporte de análisis}
  \begin{columns}
    \column{0.5\textwidth}
      \textbf{Estructura clásica (pirámide invertida):}
      \begin{enumerate}
        \item \textbf{Conclusión principal} (primero)
        \item Hallazgos que la sostienen
        \item Metodología resumida
        \item Limitaciones y próximos pasos
        \item Apéndice técnico (si hace falta)
      \end{enumerate}
      \bigskip
      El lector no técnico lee 1 y 2.\\
      El lector técnico lee todo.
    \column{0.5\textwidth}
      \begin{exampleblock}{Por qué la conclusión va primero}
        En ciencia, la hipótesis va al principio y los resultados al final.\\
        En comunicación ejecutiva, el destinatario tiene poco tiempo. Quiere saber qué hacer antes de saber cómo llegaste.\\
        \bigskip
        Adaptar la estructura a la audiencia no es menos riguroso: es mejor comunicación.
      \end{exampleblock}
  \end{columns}
\end{frame}

\begin{frame}{Narrativa: conectar los hallazgos}
  Un buen reporte no es una lista de gráficos. Es una historia con hilo conductor.

  \bigskip
  \textbf{Preguntas que guían la narrativa:}
  \begin{itemize}
    \item ¿Cuál es el problema que motivó el análisis?
    \item ¿Qué hipótesis teníamos antes de ver los datos?
    \item ¿Qué encontramos? ¿Coincidió con la hipótesis?
    \item ¿Qué implica ese hallazgo para la decisión original?
    \item ¿Qué no pudimos responder y por qué?
  \end{itemize}

  \bigskip
  \begin{alertblock}{Trampa frecuente}
    Reportar \textit{todo} lo que se analizó, aunque no sea relevante para la pregunta original. Más análisis $\neq$ mejor reporte. La selección y el énfasis son parte del trabajo.
  \end{alertblock}
\end{frame}

%-----------------------------------------------
\section{Ética en la comunicación}
%-----------------------------------------------

\begin{frame}{Honestidad e integridad en el reporte}
  \begin{columns}
    \column{0.5\textwidth}
      \textbf{Prácticas deshonestas (aunque comunes):}
      \begin{itemize}
        \item \textit{Cherry-picking}: mostrar solo los subgrupos donde el modelo funciona bien.
        \item Escala manipulada para exagerar diferencias.
        \item Omitir casos donde el modelo falla.
        \item Presentar correlación como causalidad.
        \item Redondear métricas favorablemente.
      \end{itemize}
    \column{0.5\textwidth}
      \textbf{Prácticas honestas:}
      \begin{itemize}
        \item Mostrar métricas en el conjunto de test, no de entrenamiento.
        \item Reportar intervalos de confianza.
        \item Mencionar explícitamente las limitaciones.
        \item Distinguir entre lo que los datos dicen y lo que el analista interpreta.
        \item Separar la descripción de la recomendación.
      \end{itemize}
  \end{columns}

  \bigskip
  \begin{alertblock}{Principio}
    Un análisis honesto que no respalda la hipótesis del cliente es más valioso que uno que la respalda artificialmente.
  \end{alertblock}
\end{frame}

%-----------------------------------------------
\section{Cierre}
%-----------------------------------------------

\begin{frame}{Ideas para llevarse}
  \begin{enumerate}
    \item Comunicar resultados es parte del trabajo del analista, no un extra.
    \item La audiencia determina el formato, el lenguaje y el nivel de detalle.
    \item Confirmar lo que ya se sabe tiene valor si se cuantifica y se da contexto.
    \item Un hallazgo contraintuitivo no es un error: es una oportunidad para explorar juntos.
    \item El gráfico correcto depende del mensaje, no de la estética.
    \item La conclusión va primero cuando se habla con tomadores de decisión.
    \item La honestidad sobre las limitaciones construye credibilidad a largo plazo.
  \end{enumerate}
\end{frame}

\end{document}
